% GENERATED FILE. Edit canonical lessons, then run npm run generate:books.
% canonical-source-hash: fnv1a64:d02aa229d5f81aee
% canonical-lessons: SA-C461-mahatma, SA-C461-khalah, SA-C461-saujanyam, SA-C461-nisthurah, SA-C461-druhyati

\chapter{Good People and Bad}
\label{ch:sa-good-people-and-bad}
\hlchaptermodality{461}

\begin{chapteropening}
\textbf{By the end of this chapter:} \emph{I can say whether someone is good-hearted or wicked.}

Complete the last lesson of chapter 461: I can say whether someone is good-hearted or wicked.
\end{chapteropening}

\section[\texorpdfstring{\sk{महात्मा} (mahātmā)}{महात्मा (mahātmā)}]{\sk{महात्मा} (mahātmā) — a great soul; a noble person}

\label{lesson:SA-C461-mahatma}



\begin{quote}
Before the new one: say the Sanskrit for a virtue, then the Sanskrit for wicked.
\end{quote}

\subsection*{You'll want to know: \sk{महात्मा}}
\textbf{\sk{महात्मा}} — \emph{mahātmā} — \textquotedblleft{}a great soul; a noble person\textquotedblright{}.

\begin{quote}
\sk{महात्मा} \sk{दरिद्रस्य} \sk{साहाय्यं} \sk{करोति।}
\end{quote}

\begin{quote}
\emph{mahātmā daridrasya sāhāyyaṃ karoti.}
\end{quote}

\begin{quote}
A noble soul helps the poor.
\end{quote}

A masculine noun.

\begin{cousinweb}
\textbf{\sk{महात्मा}} (\emph{mahātmā}) is built from \textbf{\sk{महत्}}, \textquotedblleft{}great\textquotedblright{}, and \textbf{\sk{आत्मा}}, \textquotedblleft{}the self, the soul\textquotedblright{}; MW gives \textquotedblleft{}high-souled, magnanimous, noble\textquotedblright{} (MW).
\end{cousinweb}

\begin{grammarlens}[title={What you've built}]
One more word for describing what you lived and felt.
\end{grammarlens}

\subsection*{Your turn}
\begin{itemize}
\raggedright
  \item \emph{Say it:} \emph{mahātmā}
  \item \emph{Say it:} \emph{mahātmā}, once more
  \item \emph{Say it:} \emph{mahātmā daridrasya sāhāyyaṃ karoti.} — A noble soul helps the poor.
  \item \emph{Recall:} say the Sanskrit for a virtue, then the Sanskrit for wicked, then say \emph{mahātmā} again
\end{itemize}

\begin{tcolorbox}[breakable,colback=teal!4,colframe=teal!35!black,title={Before you move on}]
What does \sk{महात्मा} mean? (\textquotedblleft{}A great soul; a noble person\textquotedblright{}.) Say it once more.
\end{tcolorbox}
\section[\texorpdfstring{\sk{खलः} (khalaḥ)}{खलः (khalaḥ)}]{\sk{खलः} (khalaḥ) — a scoundrel}

\label{lesson:SA-C461-khalah}



\begin{quote}
Before the new one: say the Sanskrit for wicked, then the Sanskrit for a great soul.
\end{quote}

\subsection*{You'll want to know: \sk{खलः}}
\textbf{\sk{खलः}} — \emph{khalaḥ} — \textquotedblleft{}a scoundrel\textquotedblright{}.

\begin{quote}
\sk{खलः} \sk{सज्जनं} \sk{पीडयति।}
\end{quote}

\begin{quote}
\emph{khalaḥ sajjanaṃ pīḍayati.}
\end{quote}

\begin{quote}
A scoundrel torments a good man.
\end{quote}

A masculine noun.

\begin{cousinweb}
\textbf{\sk{खलः}} (\emph{khalaḥ}) MW's first \emph{khala} is \textquotedblleft{}a threshing-floor\textquotedblright{}; a later sense is \textquotedblleft{}a mischievous man\textquotedblright{}, and for it MW gives no derivation, so its origin is uncertain (MW).
\end{cousinweb}

\begin{grammarlens}[title={What you've built}]
One more word for describing what you lived and felt.
\end{grammarlens}

\subsection*{Your turn}
\begin{itemize}
\raggedright
  \item \emph{Say it:} \emph{khalaḥ}
  \item \emph{Say it:} \emph{khalaḥ}, once more
  \item \emph{Say it:} \emph{khalaḥ sajjanaṃ pīḍayati.} — A scoundrel torments a good man.
  \item \emph{Recall:} say the Sanskrit for wicked, then the Sanskrit for a great soul, then say \emph{khalaḥ} again
\end{itemize}

\begin{tcolorbox}[breakable,colback=teal!4,colframe=teal!35!black,title={Before you move on}]
What does \sk{खलः} mean? (\textquotedblleft{}A scoundrel\textquotedblright{}.) Say it once more.
\end{tcolorbox}
\section[\texorpdfstring{\sk{सौजन्यम्} (saujanyam)}{सौजन्यम् (saujanyam)}]{\sk{सौजन्यम्} (saujanyam) — goodness; kindness}

\label{lesson:SA-C461-saujanyam}



\begin{quote}
Before the new one: say the Sanskrit for a great soul, then the Sanskrit for a scoundrel.
\end{quote}

\subsection*{You'll want to know: \sk{सौजन्यम्}}
\textbf{\sk{सौजन्यम्}} — \emph{saujanyam} — \textquotedblleft{}goodness; kindness\textquotedblright{}.

\begin{quote}
\sk{तस्य} \sk{सौजन्येन} \sk{सर्वे} \sk{प्रसन्नाः} \sk{भवन्ति।}
\end{quote}

\begin{quote}
\emph{tasya saujanyena sarve prasannāḥ bhavanti.}
\end{quote}

\begin{quote}
Everyone is pleased by his kindness.
\end{quote}

A neuter noun.

\begin{cousinweb}
\textbf{\sk{सौजन्यम्}} (\emph{saujanyam}) is built on \emph{su-jana}, \textquotedblleft{}a good person\textquotedblright{} (\emph{su}, \textquotedblleft{}good\textquotedblright{}, and \textbf{\sk{जनः}}), with the lengthening of its first vowel; MW gives \textquotedblleft{}goodness, kindness, benevolence\textquotedblright{} (MW).
\end{cousinweb}

\begin{grammarlens}[title={What you've built}]
One more word for describing what you lived and felt.
\end{grammarlens}

\subsection*{Your turn}
\begin{itemize}
\raggedright
  \item \emph{Say it:} \emph{saujanyam}
  \item \emph{Say it:} \emph{saujanyam}, once more
  \item \emph{Say it:} \emph{tasya saujanyena sarve prasannāḥ bhavanti.} — Everyone is pleased by his kindness.
  \item \emph{Recall:} say the Sanskrit for a great soul, then the Sanskrit for a scoundrel, then say \emph{saujanyam} again
\end{itemize}

\begin{tcolorbox}[breakable,colback=teal!4,colframe=teal!35!black,title={Before you move on}]
What does \sk{सौजन्यम्} mean? (\textquotedblleft{}Goodness; kindness\textquotedblright{}.) Say it once more.
\end{tcolorbox}
\section[\texorpdfstring{\sk{निष्ठुरः} (niṣṭhuraḥ)}{निष्ठुरः (niṣṭhuraḥ)}]{\sk{निष्ठुरः} (niṣṭhuraḥ) — heartless}

\label{lesson:SA-C461-nisthurah}



\begin{quote}
Before the new one: say the Sanskrit for a scoundrel, then the Sanskrit for goodness.
\end{quote}

\subsection*{You'll want to know: \sk{निष्ठुरः}}
\textbf{\sk{निष्ठुरः}} — \emph{niṣṭhuraḥ} — \textquotedblleft{}heartless\textquotedblright{}.

\begin{quote}
\sk{निष्ठुरः} \sk{स्वामी} \sk{सेवकं} \sk{ताडयति।}
\end{quote}

\begin{quote}
\emph{niṣṭhuraḥ svāmī sevakaṃ tāḍayati.}
\end{quote}

\begin{quote}
The heartless master beats his servant.
\end{quote}

\begin{cousinweb}
\textbf{\sk{निष्ठुरः}} (\emph{niṣṭhuraḥ}) MW asks, with a question mark, whether it comes from \emph{ni} and the root √sthā, \textquotedblleft{}to stand\textquotedblright{} (the root of \textbf{\sk{तिष्ठति}}), so its origin is uncertain; it gives \textquotedblleft{}hard, harsh, severe, cruel\textquotedblright{} (MW).
\end{cousinweb}

\begin{grammarlens}[title={What you've built}]
One more word for describing what you lived and felt.
\end{grammarlens}

\subsection*{Your turn}
\begin{itemize}
\raggedright
  \item \emph{Say it:} \emph{niṣṭhuraḥ}
  \item \emph{Say it:} \emph{niṣṭhuraḥ}, once more
  \item \emph{Say it:} \emph{niṣṭhuraḥ svāmī sevakaṃ tāḍayati.} — The heartless master beats his servant.
  \item \emph{Recall:} say the Sanskrit for a scoundrel, then the Sanskrit for goodness, then say \emph{niṣṭhuraḥ} again
\end{itemize}

\begin{tcolorbox}[breakable,colback=teal!4,colframe=teal!35!black,title={Before you move on}]
What does \sk{निष्ठुरः} mean? (\textquotedblleft{}Heartless\textquotedblright{}.) Say it once more.
\end{tcolorbox}
\section[\texorpdfstring{\sk{द्रुह्यति} (druhyati)}{द्रुह्यति (druhyati)}]{\sk{द्रुह्यति} (druhyati) — bears malice}

\label{lesson:SA-C461-druhyati}



\begin{quote}
Before the new one: say the Sanskrit for goodness, then the Sanskrit for heartless.
\end{quote}

\subsection*{You'll want to know: \sk{द्रुह्यति}}
\textbf{\sk{द्रुह्यति}} — \emph{druhyati} — \textquotedblleft{}bears malice\textquotedblright{}.

\begin{quote}
\sk{खलः} \sk{सज्जनाय} \sk{द्रुह्यति।}
\end{quote}

\begin{quote}
\emph{khalaḥ sajjanāya druhyati.}
\end{quote}

\begin{quote}
The scoundrel bears malice towards the good man.
\end{quote}

Past (he …): \sk{सः} \sk{मित्राय} \sk{द्रुग्धवान्}, \emph{saḥ mitrāya drugdhavān}.

\begin{cousinweb}
\textbf{\sk{द्रुह्यति}} (\emph{druhyati}) comes from the root √druh, which takes the dative of the person harmed (MW).
\end{cousinweb}

\begin{grammarlens}[title={What you've built}]
One more word for describing what you lived and felt.
\end{grammarlens}

\subsection*{Your turn}
\begin{itemize}
\raggedright
  \item \emph{Say it:} \emph{druhyati}
  \item \emph{Say it:} \emph{druhyati}, once more
  \item \emph{Say it:} \emph{khalaḥ sajjanāya druhyati.} — The scoundrel bears malice towards the good man.
  \item \emph{Recall:} say the Sanskrit for goodness, then the Sanskrit for heartless, then say \emph{druhyati} again
\end{itemize}

\begin{tcolorbox}[breakable,colback=teal!4,colframe=teal!35!black,title={Before you move on}]
What does \sk{द्रुह्यति} mean? (\textquotedblleft{}Bears malice\textquotedblright{}.) Say it once more.
\end{tcolorbox}
